% Source: source_md/intelligence_part2_combined_ja.md
\section{結論 — 意図的に開いておくもの}


本章は本論文を結ぶ。ここで行うのは、論証の追加ではなく、論文の射程の境界の明示である。本論文が何を示し、何を意図的に開いたまま残すかを述べ、三本の論理線を総括し、Part I への引き渡しを行う。

\bigskip


\subsection{非介入の表明}


本論文は、知性の動力学を記述してきた。$L_F$ の非閉包(\S{}3)、$L_E$ の創発(\S{}5)、停止構造の歴史的形態(\S{}6〜\S{}9)、ゲティア問題(\S{}10)、不安定領域(\S{}11)、構成的条項(\S{}12)。これらはすべて、知性がどのように動くかの記述である。

本論文は、知性がどう設計・運用・制度化されるべきかについては、何も述べない。実装の指針、倫理的処方、制度的提言——これらはいずれも本論文の射程外である。これは怠慢や留保ではなく、原理的な非介入である。

なぜ介入しないのか。停止構造の具体的な実装・倫理・制度は、歴史的偶然性・文化的文脈・技術的制約・正統性体系と不可分である。\S{}9 で示した通り、正統性は計算できない(\S{}9.4)。それは時間・責任・歴史・分散的受容に依存する、計算に還元されない過程である。動力学理論は、この計算不可能な過程に対して、解を与えることができない。動力学は、停止構造がどう動くかを記述できるが、特定の状況でどの停止を採用すべきかを決定できない。

これは理論の限界ではなく、理論の射程の正確な認識である。動力学の記述と、その使用法の決定は、カテゴリ的に異なる。

\bigskip


\subsection{動力学の説明と使用法の決定}


このカテゴリ的区別を明示する。

「あるものがどう動くか」の説明と、「そのものをどう使うべきか」の決定は、異なる種類の言明である。前者は動力学的記述であり、後者は規範的・実践的判断である。本論文は前者に徹し、後者には立ち入らない。

たとえば \S{}11.8 で、単一ノード・高速・高解像・内部停止の系が発散域を占めることを示した。これは動力学的記述である。この記述から「ゆえにそのような系を作るべきでない」という規範的結論は導かれない。動力学は、その系がどこに位置しどう振る舞うかを述べるが、その系を作るべきか否かは、動力学の外部にある価値判断・実践的文脈に属する。

このカテゴリ的区別を曖昧にすると、本論文は隠れた規範的主張を含むものとして誤読される。本論文の主張は、徹頭徹尾、動力学的記述である。それが実践においてどう使われるかは、本論文を読む者の責任において、本論文の外部で決定される。

\bigskip


\subsection{三本の論理線の総括}


本論文を貫いた三本の論理線を総括する。

\textbf{第一線 — 非閉包と外部停止}:
差から推論演算子 $F$ が導かれ、$F$ には内部不動点が存在しない(\S{}3、非閉包定理)。$L_F$ は単独で閉じない。沈黙世界では無限後退する(\S{}4)。ゆえに停止は外部化されねばならず、外部停止演算子 $E$ と外部参照場 $\Phi_{\text{ext}}$ が必然となる(\S{}5)。$\Phi_{\text{ext}}$ の歴史的諸形態が、アニミズム・神・宗教・社会である(\S{}6〜\S{}9)。

\textbf{第二線 — 閾値移動とサピエンス的特徴}:
相互拘束 $H_{ij}$ のパラメータが応答性閾値を決め、その閾値は本体論的範疇ではなく連続的に移動する(\S{}2.6)。閾値を内発的に動かす能力が $\mu_\theta$ であり、これは運動可能領域の拘束パラメータ $\theta$ を組み替える能力である(\S{}2.7)。$\mu_\theta$ は知性一般の条件ではなく、サピエンス的知性を特徴づけ、個体内での自律的加速をもたらす(\S{}5.4、\S{}12)。$\mu_\theta$ は誤り許容性と構造的に同値であり、閾値を動かせる知性は必然的に誤りうる(\S{}10.6)。

\textbf{第三線 — 三層動力学と創発停止}:
推論層 $L_F$ の運動は痕跡を残し、ログ層 $L_R$ を生む(\S{}4.6、\S{}5.2)。$L_R$ の再帰運動と外部相互拘束の関数として、創発停止層 $L_E$ が生じる(\S{}5.3)。$L_E$ は個体内にありながら外部依存的な停止アトラクターである。外部参照場 $\Phi_{\text{ext}}$ は個体間 $L_E$ ネットワーク $\Phi_{\text{net}}$ と同型である(\S{}5.8)。三層の協調が安定運用領域を、その崩壊が不安定領域を特徴づける(\S{}11)。

三本の線は \S{}11 の三軸相図で統合され、\S{}12 で構成的条項として整理された。知性一般は第一線・第三線によって持続し、サピエンス的知性は第二線によってそれを能動的・自律的に行う。

\bigskip


\subsection{VED の基底原理の知性層における現れ}


本論文の全体は、VED の単一の基底原理——静的範疇ではなく運動——の、知性層における展開であった。

VED は時間・空間・物質・生命・意識を、静的な実体や範疇としてではなく、差から生じる運動として再配置してきた。本論文は同じ操作を知性に対して行った。知性は、能力でも実体でも保存可能な対象でもなく、$L_F$・$L_R$・$L_E$ の同時運動として維持される動力学的状態である(\S{}12 最終条項)。

この再配置の最も先鋭な現れが、サピエンス的知性における $\mu_\theta$ であった。$\mu_\theta$ は運動可能領域そのものを運動させる能力——運動の二階の運動——である(\S{}12.4)。状態が領域内を動くだけでなく(第一階の運動)、領域の輪郭そのものが動く(第二階の運動)。サピエンス的知性は、この二階の運動を個体内で自律的に行う点において、VED の運動原理を知性層で最も深く具現する。

応答性も、停止も、閾値も、すべて運動として記述された。固定された世界の属性も、最終的な停止状態も、確定した閾値も、本論文には現れない。あるのは運動と、その運動が維持される動力学的状態だけである。これが VED の基底原理の、知性層における貫徹である。

\begin{figure}[H]
\centering
\includegraphics[width=\linewidth]{fig9_cross_domain_homology_ja.pdf}
\caption{領域横断ホモロジー。不完全閉包から創発アトラクターを経て持続的動態へ至る同一パターンが、物理・生命・意識・知性・社会に現れる。}
\label{fig:part2-cross-domain-homology}
\figurenote{中央のホモロジー核は、不完全閉包、創発アトラクター、持続的動態の三段からなる。五つの層カードは、この核の異なる実装である。}{本図は領域差を消去する同一視ではなく、層の自律性を保ったまま生成文法の共通性を示す。知性は特別だが、VED 理論ツリーの他層から孤立してはいない。}
\end{figure}

\bigskip


\subsection{残された開放性}


本論文が意図的に開いたまま残すものを明示する。これらは本論文の失敗や不足ではなく、本論文の射程の境界の外側にある問いである。

\textbf{哲学篇への引き渡し}:
意味・主観性・理解そのものの問題は、本論文では扱わなかった。本論文は知性の停止動力学を記述したが、その停止が主体にとって何を意味するか、理解とは何か、主観的経験はどう位置づくか——これらは VED 理論ツリーの哲学篇に引き渡される。本論文は知性層の動力学を閉じるが、その上位の問いは開いたままである。

\textbf{AI設計の実践問題}:
本論文の動力学的記述が、人工推論システムの設計・運用にどう適用されるかは、本論文の射程外である(\S{}13.1、\S{}13.2)。動力学は記述を与えるが、実践的指針は与えない。これは実装・倫理・制度の文脈において、本論文の外部で扱われるべき問いである。

\textbf{停止層発達度 $\Sigma$ の定量化}:
創発停止層 $L_E$ の発達度・活性度を測る量 $\Sigma$ は、概念的に存在するが本論文では量化しなかった(\S{}1.8)。$\Sigma$ の定量化は、臨床的・実証的射程に属し、後続研究へ開かれる。停止構造の硬直や崩壊(\S{}8.5、\S{}11)を定量的に扱うには、$\Sigma$ の精密化が必要だが、それは本論文の構造的記述の範囲を超える。

\textbf{停止汎関数の形式化}:
$L_E$ を Vol.2 の閉包エネルギー汎関数 $E(A)$ と平行に「停止汎関数の局所極小」として形式化することは、原理的に可能だが本論文では行わなかった(\S{}5.3)。この形式化は後続研究の課題である。

\bigskip


\subsection{Part I 微修正への引き渡し}


本論文の結果は、知性編 Part I のいくつかの記述を精密化する。Part I を改訂するための提案を以下に整理する。これらは Part I の誤りの訂正ではなく、Part II の三層動力学による精密化である。

\textbf{Part I \S{}8(個体知性の非閉包)の精密化}:
Part I \S{}8 は個体知性が閉じないことを論じた。Part II \S{}3 はこれを非閉包定理として証明し、さらに \S{}5 で「$L_F$ 単独では閉じないが、$L_R$ と $L_E$ を通じてなら外部依存的に閉じうる」と精密化した。Part I \S{}8 に、この (A)/(B) の射程の区別——固定 $\theta$ では閉じないが $\theta$ 組み替えで停止が創発する——を追加することが望ましい。

\textbf{Part I \S{}9(差分の地平線)への三層解釈の追加}:
Part I \S{}9 は知性の差分の地平線を論じた。Part II \S{}1.5 と \S{}3.6 は、これを $C_{max}$ 不到達性および内部不動点不在と同型として位置づけた。Part I \S{}9 に、差分の地平線が三層動力学($L_F$ の非閉包)としてどう現れるかの解釈を追加できる。

\textbf{Part I \S{}6(ゲティア問題)への $\mu_\theta$ 接続の補強}:
Part I \S{}6 はゲティア問題を構造的診断として扱った。Part II \S{}10.6 は、ゲティア構造の排除不能性が閾値移動能力 $\mu_\theta$ の存在と構造的に同値であることを示した。Part I \S{}6 に、この同値——誤り許容性と $\mu_\theta$ の構造的結合——を補強として追加できる。

\textbf{Part I \S{}5(Physarum・Transformer)の三層プロファイル再記述}:
Part I \S{}5 は Physarum と Transformer を知性の領域横断インスタンスとして示した。Part II \S{}5.4 は、両者を三層プロファイル(粘菌「ログ外・停止自律」、Transformer「ログ内・停止外」)として精密化し、$\mu_\theta$ の所在によってサピエンスと区別した。Part I \S{}5 に、この三層プロファイル比較を再記述として追加できる。重要なのは、これにより両者が「$L_F$ のみのインスタンス」ではなく「三層動力学の異なるプロファイル」として正しく位置づけられることである。

これらの微修正は、Part II 完成後の別作業として行われる。本論文はその提案を引き渡すに留める。

\bigskip


\subsection{知性は閉じない、しかし持続する}


本論文の帰結を、一文に集約する。

知性は完結しない。知性は閉じない。知性は最適化されない。$L_F$ は内部不動点を持たず(\S{}3)、完全閉包 $C_{max}$ は到達不能であり(\S{}1.5)、完全最適化は知性の終了である(\S{}11.6)。これらはすべて、知性が運動であることの帰結である。

しかし、知性は持続する。$L_E$ が外部依存的に停止を供給し(\S{}5)、その停止が改訂可能であることによって(\S{}5.8)、知性は閉じないまま持続する。閉じないことと持続しないことは、同じではない。知性は閉じないがゆえに運動し続け、運動し続けるがゆえに持続する。

サピエンス的知性は、これに加えて $\mu_\theta$ によって運動可能領域を自律的に組み替え、個体内で加速する。それは閉じない運動を、能動的に運用する知性である。誤りうるがゆえに改訂でき、改訂できるがゆえに持続する。

\bigskip


\subsection{理論の境界}


本論文は、VED 理論ツリーにおける知性層の動力学を閉じる。

ここで「閉じる」とは、\S{}3 が否定した意味での閉包ではない。それは、本論文の射程——知性の停止動力学の構造的記述——を完結させるという意味である。知性層の動力学について、本論文は述べるべきことを述べた。$L_F$・$L_R$・$L_E$ の三層、$\mu_\theta$ による加速、停止構造の歴史的形態、不安定領域、構成的条項。これらが、知性層の動力学の構造である。

残された問いは、別の層に属する。意味と理解は哲学篇に、実装と倫理は実践に、$\Sigma$ の定量化は後続研究に属する。本論文はこれらを開いたまま、知性層の動力学の記述をここで完結する。

理論はここで終わる。問いは終わらない。

\bigskip


\subsection{要約}


本章で述べたのは次の点である:

\begin{enumerate}
\item 本論文は知性の動力学を記述し、その設計・運用・制度化には介入しない。正統性が計算不可能である以上、動力学は実践的解を与えられない
\item 動力学の説明と使用法の決定はカテゴリ的に異なる。本論文は前者に徹する
\item 三本の論理線(非閉包と外部停止・閾値移動とサピエンス的特徴・三層動力学)が統合された。知性一般は第一線と第三線で持続し、サピエンス的知性は第二線で自律的に加速する
\item 本論文は VED の基底原理(静的範疇ではなく運動)の知性層における展開であり、$\mu_\theta$ は運動の二階の運動としてその原理を最も先鋭に具現する
\item 哲学篇・AI設計実践・$\Sigma$ 定量化・停止汎関数形式化が、開かれた問いとして残される
\item Part I の四箇所(\S{}8 非閉包・\S{}9 差分の地平線・\S{}6 ゲティア・\S{}5 Physarum/Transformer)への微修正が引き渡される
\item 知性は閉じない、しかし持続する。閉じないことと持続しないことは同じではない
\end{enumerate}


理論はここで終わる。問いは終わらない。

\bigskip
